\section{Background: Credit in Critic-Free Agent RL}
\label{sec:background}

\paragraph{Setting.}
An agent interacts with an environment for up to $T$ steps. At step $t$ the policy $\pi_\theta$ reads
a prompt built from the task, recent history and the admissible actions, and emits a response
$a_t = (a_t^1,\dots,a_t^{|a_t|})$: reasoning inside \texttt{<think>} tags followed by a command inside
\texttt{<action>} tags. The environment executes the command; the episode reward is sparse,
$r_T \in \{0, 1\}$ for success, plus a small penalty for malformed responses~\cite{feng2025gigpo}.
For each task a group of $G$ trajectories $\{\tau_i\}_{i=1}^{G}$ is sampled with the same initial state.

\paragraph{Clipped surrogate.}
All methods below share PPO's token-level objective~\cite{schulman2017ppo}
\begin{equation}
\begin{split}
\mathcal{J}(\theta) = \mathbb{E}\Big[\tfrac{1}{|a_t|}\textstyle\sum_{j}
  \min\big(&\rho_{t,j} A_{t,j},\, \mathrm{clip}(\rho_{t,j}, 1\!\pm\!\epsilon)\, A_{t,j}\big)\Big]\\
  &- \beta\, \mathrm{KL}(\pi_\theta \,\|\, \pi_{\mathrm{ref}}),
\end{split}
\label{eq:ppo}
\end{equation}
with $\rho_{t,j}$ the token probability ratio to the sampling policy. They differ only in the
advantage $A_{t,j}$, and until now $A_{t,j}$ has not depended on $j$.

\paragraph{GRPO: episode credit.}
GRPO~\cite{shao2024deepseekmath} standardizes each trajectory's return within its group,
$\Aep_i = (R(\tau_i) - \mu_G)/(\sigma_G + \varepsilon)$, and broadcasts it to every token of every step.

\paragraph{GiGPO: step credit from anchor states.}
GiGPO~\cite{feng2025gigpo} collects every visit, across the group's rollouts, to the same environment
state $z$ (identical observation), forming an anchor group $\mathcal{G}_z$ of size $m_z$. With the
discounted step return $\hat g_t = \sum_{t' \geq t} \gamma^{t'-t} r_{t'}$, the step advantage
standardizes $\hat g_t$ within the anchor group,
$\Ast_t = (\hat g_t - \mu_{\mathcal{G}_z})/(\sigma_{\mathcal{G}_z} + \varepsilon)$
(falling back to prompt-group statistics when $m_z = 1$~\cite{liang2026gaca}), and the two levels are
summed, $A_t = \Aep + w\,\Ast_t$.

\paragraph{GACA: state-adaptive mixing.}
GACA~\cite{liang2026gaca} replaces the sum by a per-step convex mixture
\begin{equation}
A_t = \lambda_t\, w\, \Ast_t + (1 - \lambda_t)\, \Aep.
\label{eq:gaca}
\end{equation}
The weight comes from a criticality score: the mean token NLL of the sampled response,
$H_t = -\frac{1}{|a_t|}\sum_j \log \pi_{\theta_{\mathrm{old}}}(a_t^j \mid s_t, a_t^{<j})$, normalized by its
trajectory mean, $c_t = H_t / \bar H$. A piecewise-linear gate maps it to the weight,
$\lambda_t = \lambda_0(1 + \alpha(c_t - 1))$ for $c_t \geq 1$ and $\lambda_0(1 - \beta(1 - c_t))$
otherwise, clamped to $[0, 1]$, with $\lambda_0 = 0.5$ and $\alpha = \beta = 0.5$. Uncertain steps
lean on the step contrast, confident steps on the episode return. GACA's analysis gives the
MSE-optimal weight as $\lambda^\star_z = M_E/(M_E + M_S)$, the ratio of the two estimators' errors, and
shows it is non-decreasing in $m_z$; the NLL enters as a one-sided certificate of local decision
criticality. In every case above, all $|a_t|$ tokens of a response share $A_t$.
